\textbf{Model.} With $w_{y_m}=q+b_y$ the measured pitch rate is a known input, and the bias becomes a state
modelled as a random walk ($\dot b_y=0$ plus a small $w_b$ used only for tuning):
\begin{equation*}
\underbrace{\begin{bmatrix}\dot\theta\\ \dot b_y\end{bmatrix}}_{\dot x}=
\underbrace{\begin{bmatrix}0&-1\\0&0\end{bmatrix}}_{A}x+\begin{bmatrix}1\\0\end{bmatrix}w_{y_m}+\begin{bmatrix}w\\ w_b\end{bmatrix},\qquad
y=\underbrace{\begin{bmatrix}1&0\end{bmatrix}}_{C}x+v,
\end{equation*}
which is observable (the pair $(A,C)$ has $\mathrm{rank}[C;CA]=2$). Discretised at the real packet period
$T_s=13.2$\,ms (75.6\,Hz, mean of the fresh-packet intervals, not the 200\,Hz log rate; only \texttt{nd.new} samples are used):
$x_{k+1}=\left[\begin{smallmatrix}1&-T_s\\0&1\end{smallmatrix}\right]x_k+\left[\begin{smallmatrix}T_s\\0\end{smallmatrix}\right]w_{y_m,k}+\xi_k$,
$\mathrm{cov}\,\xi=\mathrm{diag}\big((\sigma_qT_s)^2,\ \sigma_b^2T_s\big)$, $\mathrm{cov}\,v=\sigma_y^2$.
The steady-state (current-estimator) gain $L=[L_\theta\ L_b]^\top$ is computed with \texttt{kalman()} (Control System Toolbox;
cross-checked with the Riccati equation \texttt{idare}, identical to 5 digits) and applied as
$\hat x_k^-=F\hat x_{k-1}+[T_s;0]\,w_{y_m,k-1}$, $\hat x_k=\hat x_k^-+L\,(y_k-\hat\theta_k^-)$, with the actual packet interval in~$F$.
The no-bias filter of 4.1 was re-implemented in the script (same $\sigma_q,\sigma_y$) for comparison; it gives $L=0.0150$ ($k\simeq1.15$\,s$^{-1}$).

\textbf{Tuning.} $\sigma_q=0.03$\,rad/s (std of $q$ in hover, Exp.~C; at rest it is only $0.75\times10^{-3}$) and
$\sigma_y=1.5^\circ$ (rest: $0.41^\circ$; larger in flight due to vibration/specific force) are kept from 4.1. The only new knob is the bias
random-walk intensity $\sigma_b$ [rad/s$/\sqrt{\text{s}}$] (ratio $Q_b/R$): it sets how fast $\hat b_y$ may move.
\emph{Simulation} (Fig.~\ref{fig:q44sim}): an artificial bias $b_y=+0.05$\,rad/s (stepping to $+0.08$ at 30\,s) is added to a synthetic
gyro with $\theta(t)=5^\circ[\sin(2\pi0.15t)+0.6\sin(2\pi0.4t)]$ and the same noise levels. Small $\sigma_b$ gives a smooth but slow $\hat b_y$
($10^{-4}$: not converged after 60\,s; $3\times10^{-4}$: 10\,\% error reached after 25\,s, $\theta$ rms error $0.78^\circ$); large $\sigma_b$ converges
in 2\,s but $\hat b_y$ becomes noisy and follows inclinometer errors. $\sigma_b=10^{-3}$ is a compromise: 9\,s to reach 10\,\% of the bias, $\theta$ error
$0.33^\circ$ rms (raw inclinometer: $1.5^\circ$). The resulting gains are $L=[0.0184,\,-0.00435]$, i.e.\ in continuous time
$k_1=1.41$\,s$^{-1}$, $k_2=0.33$\,s$^{-2}$ ($\omega_n=0.58$\,rad/s, $\zeta=1.2$) and a bias-free filter with the same measurements has mean $\theta$ error
$3.2^\circ$ ($\approx b_y/k$), whereas the bias filter removes it (rms error $0.33^\circ$).
\begin{center}
\includegraphics[width=0.78\linewidth]{figs/q4_4_sim_bias.pdf}\\[-1mm]
\captionof{figure}{Simulation with artificial gyro bias: true $b_y$ vs.\ $\hat b_y$ for several $\sigma_b$ (top), $\theta$ error of the no-bias and bias
filters ($\sigma_b=10^{-3}$, middle), and bias estimation error (bottom).}\label{fig:q44sim}
\end{center}

\textbf{Real data} ($\sigma_b=10^{-3}$, $\hat b_y(0)=0$, $\hat\theta(0)=\theta_m(0)$). \emph{Exp.~A} (at rest): the measured mean of the pitch gyro is $-1.57\times10^{-3}$\,rad/s
and $\hat b_y$ gets within 1\,mrad/s of it after 3.5\,s and then fluctuates around it (last 10\,s: mean $-0.7$\,mrad/s, std $1.3$\,mrad/s;
Fig.~\ref{fig:q44real}a). The residual error is of the order of the bias itself because $\hat b_y$ is driven by inclinometer noise (tuned for $1.5^\circ$,
actual $0.41^\circ$ at rest) and is irrelevant for $\hat\theta$ (error $\lesssim 1\,\text{mrad/s}/k\approx0.05^\circ$).
\emph{Exp.~C} (hover): the gyro has essentially no bias ($+0.45$\,mrad/s mean), and $\hat b_y$ stays within $\pm1$\,mrad/s; the bias filter behaves like the
no-bias one (rms difference to the onboard $\theta$: raw $0.73^\circ$, KF without bias $0.52^\circ$, with bias $0.55^\circ$).
\emph{Exp.~D} (pitch pulse, 5--18.4\,s, Fig.~\ref{fig:q44real}b): both filters follow the onboard estimate with lag and reduced amplitude (the peak is $\approx4^\circ$
vs.\ $5.8^\circ$ onboard), because during the manoeuvre the accelerometer measures specific force (rotor drag, acceleration) and not just gravity, so
the inclinometer disagrees with the integrated gyro. The bias filter interprets this disagreement as bias: $\hat b_y$ swings to $+18$\,mrad/s and $-24$\,mrad/s
(std $13$\,mrad/s, versus the real bias of $\approx0.6$\,mrad/s), and the rms difference to the onboard $\theta$ rises to $1.8^\circ$ (no bias: $1.29^\circ$; raw: $3.06^\circ$).
Larger $\sigma_b$ makes this worse ($3\times10^{-3}$: $2.55^\circ$; $10^{-2}$: $3.15^\circ$), smaller is better ($3\times10^{-4}$: $1.44^\circ$, at the cost of 25\,s to converge on a large bias).
Hence the bias state is an \emph{insurance}: with this drone/gyro (bias $\ll$ noise) it brings no accuracy gain over 4.1, but it is essential when a bias such as the
artificial $0.05$\,rad/s appears, and its tuning must be conservative ($\omega_n\approx0.6$\,rad/s) so that the specific-force error of the inclinometer is not absorbed into $\hat b_y$.
\begin{center}
\includegraphics[width=0.49\linewidth]{figs/q4_4_expA_bias.pdf}\hfill\includegraphics[width=0.49\linewidth]{figs/q4_4_expC.pdf}\\[1mm]
\includegraphics[width=0.78\linewidth]{figs/q4_4_expD.pdf}\\[-1mm]
\captionof{figure}{Real data. (a) Exp.~A: $\hat b_y$ vs.\ mean of the pitch gyro (dashed). (b) Exp.~C: raw, onboard, KF without and with bias, and $\hat b_y$.
(c) Exp.~D (5--18.4\,s): same signals.}\label{fig:q44real}
\end{center}
